\startsection[title={Fermentação Alcoólica}]
	%ethanol pathways
	Biocombustivel é uma alternativa energética extensamente adotada pela industria de transporte em todo o mundo como substituto ao combustivel fossil (INSERIR DADOS DO BASIL), dentre eles, o etanol desempenha um papel expessivo na demanda energética (DADOS DE CONTRIBUIAÇAO DO ETANOL). A maior parte de todo o biocombustivel produzido e distribuido comercialmente partem de biomassa lignocelulosica, ou seja, existe uma certa dificuldade técnica de pré tratamento desse material para que o biocombustivel seja sintetizado.

	Existem duas rotas principais para produção de biocombustivel a partir de biomassa lignocelulosica, estas são bioquímicas e termoquímicas. Em rotas bioquimicas o pré tratamento é feito em meio ácido ou acalino, ou tambem com explosão a vapor visando a quebra de ligações celolose-hemicelulose-lignina \cite[authoryear][ref::munasinghe_0]. Essa quebra permite a ação enzimatica de hidrólise, culminando na formação de açucares fermentáveis. Em contrapartida, o tratamento da biomassa nas rotas termoquímicas envolve a gaseificação da biomassa em gás de sintese (\chemical{CO} e \chemical{H_{2}}), para posteriormente ser convertido em biocombustivel, usando catalisadores químicos no processo de {\emph Fischer-Tropsch}, ou microorganismos catalisadores no processo conhecido como fermentação de gás de sintese ou {\emph syngas fermentation} \cite[authoryear][ref::munasinghe_0].

	No que se refere à produção de álcool a partir de fontes imediatas de carbono, como glicose ou sacarose, as leveduras desempenham papel fundamental. Em especial, \getbuffer[sacc] representa um grupo extremamente relevante nesse processo. Do ponto de vista evolutivo, essas leveduras desenvolveram essa capacidade inicialmente em resposta ao surgimento de frutas, fontes naturais de açúcares simples, possibilitando a conversão desses açúcares em álcool tanto em condições aeróbicas quanto anaeróbicas \cite[authoryear][ref::dashko_0].

	Em geral, o metabolismo de um microrganismo pode ser modificado por meio da inserção de genes reacionais em estequiometria adequada para maximizar determinado produto ou inibir mecanismos metabólicos específicos. \cite[authoryears][ref::meadows_0] trabalharam na modificação do genoma de variantes de \getbuffer[sacc] e demonstrou a possibilidade de aumentar significativamente a seletividade de caminhos metabólicos de interesse, especificamente na produção de precursores de ácidos isoprenoides utilizados no tratamento da malária. Em situações como essa, estudadas por \cite[author][ref::meadows_0], o metabolismo torna-se diferente daquele naturalmente encontrado em variantes de \getbuffer[sacc]. Entretanto, ao restringir a análise ao metabolismo nativo genérico, caracterizado principalmente pelo consumo de glicose e pela formação de etanol, é possível visualizar o mecanismo representado na \in{Figura}[fig:alcoholic_fermentation].

	\startplacefigure[title={Mecanismo genérico parcial de fermentação alcoólica}, reference={fig:alcoholic_fermentation}]
		\externalfigure[image/alcoholic_fermentation.png][height=200pt]
	\stopplacefigure

	A \in{Figura}[fig:alcoholic_fermentation] mostra a esquematização da rota de produção de etanol na fermentação alcoólica, onde os açúcares, especialmente glicose e sacarose, são metabolizados pelas leveduras em condições predominantemente anaeróbicas. Nesse processo, a glicose é convertida em piruvato por glicólise e, posteriormente, em etanol e dióxido de carbono. Sendo que durante a glicólise há formação de \chemical{ATP} e redução de \chemical{NAD^{+}}, em que este é responsavel pelo transporte de elétrons no ciclo produtivo, enquanto que na formação de etanol a partir de \chemical{Acetaldeido} o \chemical{NADH} é oxidado, ou seja, um ciclo de oxidação-redução do \chemical{NADH}.

	\startplaceformula[reference={eq:alcoholic_fermentation}]
		\startformula
			\startchemicalformula
				\chemical{C_{6}H_{12}O_{6}}
				\chemical{GIVES}
				\chemical{{2}CO_{2}}
				\chemical{PLUS}
				\chemical{{2}C_{2}H_{5}OH}
			\stopchemicalformula
		\stopformula
	\stopplaceformula

	A \in{Equação}[eq:alcoholic_fermentation] coloca a \in{Figura}[fig:alcoholic_fermentation] em uma estequimetria global do processo, sem levar em consideração os ciclos de renegeneração de \chemical{NADH} e \chemical{ATP}. Alem de que ambas representações deixam de fora considerações sbre mecanismos alternativos recorrentes como a produção de {\bf burirato}, ácido acético e butanol e intermediarios equivalentes a \chemical{Acetyl-CoA}. Na \in{Equação}[eq:alcoholic_fermentation], vê-se formação de dióxido de carbono e etanol deterimento do consumo de glicose em uma rota preferêncial de processos produtivos que visam a maximização de conversão em etanol.
\stopsection

